%% ==========================================================================
%%  SECTION 7 --- HOW THIS LITERATURE IS EVALUATED               [Track T4]
%%
%%  Job: two purposes at once. (a) Give the reader the tools to judge the
%%  numbers in \S4 and \S5. (b) Commit US, in writing, to the same standard --
%%  so the final report is measured against a bar we set ourselves.
%%  ~1 page. Short, concrete, and unusually valuable: most surveys skip it.
%% ==========================================================================
\section{How This Literature Is Evaluated}
\label{sec:eval}

\todo{Open by justifying the section's existence: the numbers in \S4 and \S5
      are not comparable unless the reader knows what was held fixed. This
      section makes the comparison rules explicit --- and commits us to them.}

\subsection{Metrics that matter}
\label{sec:eval:metrics}

\todo{Four families, each with an example from the surveyed work:
      (a) wall-clock --- but always paired with hardware. \nudgeplain{}'s
          50 minutes is on THREE 192-CORE MACHINES; the number is meaningless
          without that.
      (b) communication --- server-to-server bytes (40\,GB for \nudgeplain{} on
          Netflix) and per-client bytes (298\,KB serving).
      (c) client-side cost --- storage and bandwidth. This is what separates
          Pacmann (614\,MB client storage) from Tiptoe.
      (d) QUALITY --- nDCG@$\topk$, Recall@$\topk$. A private recommender that
          recommends badly has solved nothing. \nudgeplain{} scores 0.29 nDCG@20
          against 0.31 for non-private neural models; that gap IS the result.}

\subsection{Network models}
\label{sec:eval:network}

\todo{Make the point sharply, because it is where course projects fail:
      results in this literature FLIP under network constraints, because
      round complexity starts to dominate bandwidth. A localhost number tells
      you almost nothing about a deployment.
      Standard practice is \texttt{tc netem} with LAN and WAN profiles ---
      e.g.\ 30\,ms/100\,Mbps and 100\,ms/10\,Mbps. Note which surveyed papers
      report WAN numbers and which report LAN only. \nudgeplain{}'s headline is
      LAN. \verify{check each paper's setup section before asserting this}}

\subsection{Baselines}
\label{sec:eval:baselines}

\todo{The three baselines this literature expects, and why each is needed:
      (a) the plaintext system --- gives the honest slowdown factor;
      (b) a generic-MPC-framework implementation (MP-SPDZ) --- the ``why not
          just use a framework'' answer;
      (c) at least one published secure system --- the credibility test.
      Observe that the strongest papers report the UNFLATTERING comparison
      too. A slowdown stated plainly reads as confidence; a slowdown omitted
      reads as evasion.}

\subsection{The standard we hold ourselves to}
\label{sec:eval:ourstandard}

\todo{Write this as a commitment, in the first person, and keep it to five
      bullets. It is quoted back at us in the final report:
      (1) a written threat model and an explicit leakage profile;
      (2) baselines including the unflattering ones;
      (3) WAN numbers, not just localhost;
      (4) a microbenchmark breakdown --- where do the milliseconds go;
      (5) reproducibility: one command per figure, a README a stranger can
          follow.
      None of these requires cryptographic novelty, and most groups will do
      none of them.}
